\chapter{Kod źródłowy mostka vCAN (bridge\_client.py)}
\label{app:bridge}

Poniżej przedstawiono pełny kod pliku \texttt{bridge\_client.py} (wersja dwukierunkowa z obsługą auto-reconnect, filtrowania ID oraz wysyłania ramek z vcan0 do serwera).

\lstset{language=Python, basicstyle=\ttfamily\footnotesize, breaklines=true, frame=single, numbers=left, numberstyle=\tiny, caption={}, captionpos=b}
\begin{lstlisting}
\#!/usr/bin/env python3
"""
Mostek WebSocket -> wirtualny CAN (vcan0) z auto-reconnect i filtrowaniem ID.
Dwukierunkowy: odbiera ramki z WebSocket -> vcan0 oraz wysyła ramki z vcan0 -> WebSocket.
"""

import json
import logging
import threading
import time
import subprocess
import sys
import argparse
from typing import List, Optional

import websocket

try:
    import can
    from can import ThreadSafeBus
except ImportError:
    can = None
    ThreadSafeBus = None

logging.basicConfig(level=logging.INFO, format='\%(asctime)s [\%(levelname)s] \%(message)s')
logger = logging.getLogger("CAN-Bridge")


class VCanBridge:
    def \_\_init\_\_(self, url: str, token: Optional[str] = None,
                 filter\_ids: Optional[List[int]] = None,
                 auto\_reconnect: bool = True,
                 reconnect\_delay: int = 5,
                 status\_callback=None):
        self.url = url
        self.token = token
        self.filter\_ids = set(filter\_ids) if filter\_ids else None
        self.auto\_reconnect = auto\_reconnect
        self.reconnect\_delay = reconnect\_delay
        self.status\_callback = status\_callback

        self.ws = None
        self.bus = None
        self.running = False
        self.reconnect\_thread = None

        self.sending\_to\_vcan = False   \# flaga zapobiegająca pętli zwrotnej
        self.tx\_thread = None

        self.\_setup\_vcan()

    def \_setup\_vcan(self):
        if can is None:
            raise RuntimeError("python-can not installed; bridge cannot run.")
        try:
            subprocess.run(['ip', 'link', 'show', 'vcan0'], check=True, capture\_output=True)
            logger.info("vcan0 już istnieje.")
        except subprocess.CalledProcessError:
            logger.info("Tworzenie vcan0...")
            subprocess.run(['sudo', 'ip', 'link', 'add', 'dev', 'vcan0', 'type', 'vcan'], check=True)
            subprocess.run(['sudo', 'ip', 'link', 'set', 'vcan0', 'up'], check=True)
            logger.info("vcan0 utworzony i włączony.")
        self.bus = ThreadSafeBus(channel='vcan0', interface='socketcan')

    def \_should\_forward(self, can\_id: int) -> bool:
        if self.filter\_ids is None:
            return True
        return can\_id in self.filter\_ids

    def \_on\_message(self, ws, message):
        try:
            data = json.loads(message)
            can\_id = int(data['id'], 16) if isinstance(data['id'], str) else data['id']

            if not self.\_should\_forward(can\_id):
                logger.debug(f"Pominięto ID 0x\{can\_id:X\} (filtrowanie)")
                return

            payload = bytes(data['data'])
            is\_extended = data.get('is\_extended', False)
            msg = can.Message(arbitration\_id=can\_id, data=payload, is\_extended\_id=is\_extended)
            self.sending\_to\_vcan = True
            self.bus.send(msg)
            self.sending\_to\_vcan = False
            if self.status\_callback:
                self.status\_callback(f"Odebrano z serwera: \{msg\}")
        except Exception as e:
            logger.error(f"Błąd ramki: \{e\}")

    def \_on\_error(self, ws, error):
        logger.error(f"WebSocket error: \{error\}")
        if self.status\_callback:
            self.status\_callback(f"Błąd: \{error\}")

    def \_on\_close(self, ws, code, msg):
        logger.info("WebSocket zamknięty.")
        if self.status\_callback:
            self.status\_callback("Rozłączono")
        if self.auto\_reconnect and self.running:
            self.\_start\_reconnect()

    def \_on\_open(self, ws):
        logger.info("Połączono z serwerem.")
        if self.token:
            ws.send(self.token)
        \# Wyślij listę dozwolonych ID (jeśli ustawiono)
        if self.filter\_ids is not None:
            allowed\_msg = json.dumps(\{"allowed\_ids": list(self.filter\_ids)\})
            ws.send(allowed\_msg)
            logger.info(f"Wysłano listę dozwolonych ID: \{self.filter\_ids\}")

        if self.status\_callback:
            self.status\_callback("Połączono")

    def \_start\_reconnect(self):
        if self.reconnect\_thread and self.reconnect\_thread.is\_alive():
            return
        self.reconnect\_thread = threading.Thread(target=self.\_reconnect\_loop, daemon=True)
        self.reconnect\_thread.start()

    def \_reconnect\_loop(self):
        while self.running:
            logger.info(f"Próba ponownego połączenia za \{self.reconnect\_delay\} s...")
            time.sleep(self.reconnect\_delay)
            if not self.running:
                break
            try:
                self.ws = websocket.WebSocketApp(
                    self.url,
                    on\_open=self.\_on\_open,
                    on\_message=self.\_on\_message,
                    on\_error=self.\_on\_error,
                    on\_close=self.\_on\_close
                )
                wst = threading.Thread(target=self.ws.run\_forever, daemon=True)
                wst.start()
                break
            except Exception as e:
                logger.error(f"Błąd reconnect: \{e\}")

    def start(self):
        self.running = True
        self.ws = websocket.WebSocketApp(
            self.url,
            on\_open=self.\_on\_open,
            on\_message=self.\_on\_message,
            on\_error=self.\_on\_error,
            on\_close=self.\_on\_close
        )
        wst = threading.Thread(target=self.ws.run\_forever, daemon=True)
        wst.start()

        \# Uruchom wątek nadawczy (vcan -> WebSocket)
        self.tx\_thread = threading.Thread(target=self.\_tx\_loop, daemon=True)
        self.tx\_thread.start()

    def \_tx\_loop(self):
        """Wątek odbierający ramki z vcan0 i wysyłający je przez WebSocket."""
        while self.running:
            try:
                msg = self.bus.recv(timeout=0.5)
                if msg is None:
                    continue
                if self.sending\_to\_vcan:
                    continue
                if self.filter\_ids and msg.arbitration\_id not in self.filter\_ids:
                    continue
                if self.ws and self.ws.sock and self.ws.sock.connected:
                    frame = \{
                        "id": hex(msg.arbitration\_id),
                        "data": list(msg.data),
                        "is\_extended": msg.is\_extended\_id
                    \}
                    self.ws.send(json.dumps(frame))
                    if self.status\_callback:
                        self.status\_callback(f"Wysłano do serwera: \{msg\}")
            except Exception as e:
                logger.error(f"Błąd w pętli TX: \{e\}")
                time.sleep(0.1)

    def stop(self):
        self.running = False
        if self.tx\_thread:
            self.tx\_thread.join(timeout=1.0)
        if self.ws:
            self.ws.close()
        if self.bus:
            self.bus.shutdown()
        if self.status\_callback:
            self.status\_callback("Zatrzymano")


def main():
    parser = argparse.ArgumentParser()
    parser.add\_argument('--url', default='ws://localhost:8765')
    parser.add\_argument('--token')
    parser.add\_argument('--filter-ids', help='Lista ID CAN oddzielonych przecinkami, np. 0x123,0x456')
    parser.add\_argument('--no-auto-reconnect', action='store\_true', help='Wyłącz automatyczne ponawianie')
    args = parser.parse\_args()

    filter\_ids = None
    if args.filter\_ids:
        try:
            filter\_ids = [int(x.strip(), 16) if x.strip().startswith('0x') else int(x.strip())
                          for x in args.filter\_ids.split(',')]
        except ValueError:
            logger.error("Nieprawidłowy format listy ID.")
            sys.exit(1)

    bridge = VCanBridge(
        url=args.url,
        token=args.token,
        filter\_ids=filter\_ids,
        auto\_reconnect=not args.no\_auto\_reconnect,
        status\_callback=print
    )
    bridge.start()
    try:
        while True:
            time.sleep(1)
    except KeyboardInterrupt:
        bridge.stop()


if \_\_name\_\_ == "\_\_main\_\_":
    main()
\end{lstlisting}
